%% ============================================================
%% Episode 08 - Farsi Podcast Script (Translation)
%% Source: specialized_advisor_report_complete/markdown/08_boundary_layer_mt6.md
%% Series: Advisor Progress Report - Deep Dive
%% Compiler: XeLaTeX ONLY
%% Font: B Nazanin
%% CRITICAL: xepersian must always be loaded LAST
%% ============================================================

\documentclass[12pt, a4paper]{article}

%% ============================================================
%% COMPATIBILITY SHIMS
%% ============================================================
\makeatletter
\providecommand\IfClassLoadedT[2]{\@ifclassloaded{#1}{#2}{}}
\providecommand\IfClassLoadedF[2]{\@ifclassloaded{#1}{}{#2}}
\providecommand\IfPackageLoadedT[2]{\@ifpackageloaded{#1}{#2}{}}
\providecommand\IfPackageLoadedF[2]{\@ifpackageloaded{#1}{}{#2}}
\providecommand\IfFileLoadedTF[3]{\@ifpackageloaded{#1}{#2}{#3}}
\providecommand\IfFileLoadedT[2]{\@ifpackageloaded{#1}{#2}{}}
\providecommand\IfFileLoadedF[2]{\@ifpackageloaded{#1}{}{#2}}
\makeatother
\usepackage{expl3}
\ExplSyntaxOn
\cs_if_exist:NF \prop_gput_if_not_in:Nnn
  {
    \cs_new:Npn \prop_gput_if_not_in:Nnn #1#2#3
      { \prop_if_in:NnF #1 {#2} { \prop_gput:Nnn #1 {#2} {#3} } }
  }
\ExplSyntaxOff
%% ============================================================

\usepackage[
  a4paper,
  top=2.5cm, bottom=2.5cm,
  left=2.8cm, right=2.8cm,
  headheight=25pt
]{geometry}

\usepackage{xcolor}
\definecolor{seriesblue}{RGB}{25, 80, 150}
\definecolor{audionote}{RGB}{130, 50, 15}
\definecolor{sectiongray}{RGB}{75, 75, 75}
\definecolor{audiotan}{RGB}{252, 243, 233}

\usepackage{amsmath}
\usepackage{amssymb}
\usepackage{enumitem}
\usepackage{setspace}
\setstretch{1.8}
\usepackage{mdframed}
\usepackage{titlesec}

%% ============================================================
%% xepersian MUST be loaded LAST
%% ============================================================
\usepackage{xepersian}
\settextfont[Scale=1.05]{B Nazanin}
\setlatintextfont{Times New Roman}

\newcommand{\dash}{\lr{---}}

%% ============================================================
%% Page style -- savebox approach
%% ============================================================
\newsavebox{\hdrLeft}
\newsavebox{\hdrRight}
\AtBeginDocument{%
  \savebox{\hdrLeft}{\normalfont\small\color{sectiongray}%
    گزارش پيشرفت مشاور / بررسی عمیق}%
  \savebox{\hdrRight}{\normalfont\small\color{sectiongray}%
    قسمت هشتم: مطالعه لايه مرزی}%
}
\makeatletter
\def\ps@podcast{%
  \def\@oddhead{\usebox{\hdrLeft}\hfill\usebox{\hdrRight}}%
  \let\@evenhead\@oddhead
  \def\@oddfoot{\normalfont\hfil\thepage\hfil\relax}%
  \let\@evenfoot\@oddfoot
  \let\@mkboth\markboth
}
\makeatother
\pagestyle{podcast}

%% ============================================================
%% Section formatting
%% ============================================================
\titleformat{\section}
  {\large\bfseries\color{seriesblue}}
  {}
  {0pt}
  {}
  [\vspace{-2pt}{\color{seriesblue}\rule{\linewidth}{0.7pt}}\vspace{2pt}]

\titleformat{\subsection}
  {\normalsize\bfseries\color{sectiongray}}
  {}
  {0pt}
  {}

\setcounter{secnumdepth}{0}

\setlist[itemize]{
  itemsep=4pt,
  topsep=5pt,
  parsep=0pt,
  label=$\bullet$
}

%% ============================================================
\begin{document}
%% ============================================================

\thispagestyle{empty}

\vspace*{1.2cm}

{\centering
{\color{seriesblue}\fontsize{24}{30}\selectfont\bfseries قسمت هشتم}

\vspace{0.5cm}
{\fontsize{15}{20}\selectfont\bfseries مطالعه لايه مرزی: وقتی نتيجه اول اشتباه بود}

\vspace{0.9cm}
{\color{seriesblue}\rule{0.55\linewidth}{1pt}}
\vspace{0.7cm}

{\normalsize
\begin{tabular}{rl}
  \textbf{مجموعه:}     & گزارش پيشرفت مشاور \dash{} بررسی عمیق \\[5pt]
  \textbf{مدت:}        & ۸ تا ۱۰ دقیقه \\[5pt]
  \textbf{راوی:}       & مجری تنها \\[5pt]
  \textbf{منابع:}      & بخش ۶ گزارش مشاور \\
\end{tabular}
}

\vspace{0.9cm}
{\color{seriesblue}\rule{0.55\linewidth}{0.4pt}}
\par}

\vspace{1.2cm}

%% ============================================================
\section{افتتاحیه: ادعايی که بايد تغيير می‌کرد}

مطالعه لايه مرزی \lr{MT-6} در ابتدا ادعا کرد کاهش ۶۶.۵ درصدی در لرزش وقتی پارامتر لايه مرزی اپسيلون بهينه شد. آن عدد اشتباه است. رقم اصلاح‌شده ۳.۷ درصد است.

اين يک اختلاف گرد کردن نيست. ضريب ۱۸ است. درک اينکه چرا عدد اصلی دست‌بالا داده شد \dash{} و چرا اصلاح قابل اعتماد است \dash{} قبل از مواجهه با مشاور ضروری است.

%% ============================================================
\section{پارامتر لايه مرزی چه کاری می‌کند؟}

ابتدا يک يادآوری کوتاه. در کنترل مُد لغزشی کلاسيک، قانون کنترل شامل جمله \lr{K} ضربدر اشباع سيگما بر اپسيلون است. تابع اشباع تکه‌تکه است: داخل يک باند به عرض اپسيلون دور سطح لغزشی، کنترل خطی تغيير می‌کند. بيرون از آن باند، به شکل سخت سوئيچ می‌کند.

مبادله مستقيم است. اپسيلون بزرگ يعنی يک ناحيه نرم گسترده \dash{} کنترل ملايم و لرزش پايين، اما کنترل‌گر يک باند تحمل بزرگ‌تر دور سطح قبول می‌کند. اپسيلون کوچک يعنی ناحيه نرم باريک \dash{} کنترل تهاجمی‌تر که سطح را دقيق‌تر دنبال می‌کند، اما سوئيچينگ فرکانس بالا افزايش می‌يابد.

مطالعه \lr{MT-6} پرسيد: کدام مقدار اپسيلون لرزش را به حداقل می‌رساند در حالی که دقت رديابی قابل قبولی حفظ می‌کند؟ سه مقدار آزمايش شد: ۰.۰۱، ۰.۰۲، و ۰.۰۵.

%% ============================================================
\section{معيار اصلی و نقص آن}

اندازه‌گيری اصلی لرزش يک محاسبه \lr{RMS} حوزه زمان بود: سيگنال کنترل \lr{u} را در شبيه‌سازی ۱۰ ثانيه کامل بگيريد، ميانگين مجذور ريشه را محاسبه کنيد. مقادير \lr{RMS} را در تنظيمات مختلف اپسيلون مقايسه کنيد.

مشکل: در يک اجرای پايدارسازی معمول، ۲ تا ۳ ثانيه اول فاز دسترسی است \dash{} کنترل‌گر سخت کار می‌کند تا سيستم را از شرط اوليه به سطح لغزشی برساند. در طول اين فاز، سيگنال کنترل صرف نظر از اپسيلون بزرگ است و سريع تغيير می‌کند. \lr{RMS} در شبيه‌سازی کامل تحت سلطه اين گذرا است.

وقتی اپسيلون از ۰.۳ به ۰.۰۲ تغيير می‌کند، ديناميک‌های فاز دسترسی به شکل قابل توجهی تغيير می‌کنند \dash{} اپسيلون کوچک‌تر يعنی يک فشار اوليه تهاجمی‌تر، که يک گذرای بزرگ‌تر توليد می‌کند. معيار سيب‌های مختلف را با هم مقايسه می‌کرد: عمدتاً گذرای دسترسی را اندازه‌گيری می‌کرد، نه لرزش حالت ماندگار که طراحی لايه مرزی در واقع برای آن است.

عدد اصلی ۶۶.۵ درصد اين مصنوع را منعکس می‌کرد.

%% ============================================================
\section{روش اصلاح‌شده}

آناليز اصلاح‌شده به جای حوزه فرکانسی استفاده می‌کند. به طور خاص: سيگنال کنترل را بعد از اينکه سيستم به حالت ماندگار رسيد بگيريد \dash{} بعد از اينکه گذرای اوليه خاموش شد \dash{} يک تبديل فوريه سريع اعمال کنيد، و انرژی در باند فرکانس بالا بالای ۲۰ هرتز را اندازه‌گيری کنيد.

اين مستقيماً لرزش را اندازه‌گيری می‌کند. لرزش به صورت نوسان فرکانس بالا در سيگنال کنترل ظاهر می‌شود. \lr{FFT} محتوای فرکانسی را به صورت پاک جدا می‌کند. با پنجره‌بندی ثانيه‌های اول گذرا و تمرکز بر رفتار حالت ماندگار، معيار آنچه اپسيلون واقعاً طراحی شده تأثير بگذارد را اندازه‌گيری می‌کند.

نتيجه اصلاح‌شده: کاهش لرزش از اپسيلون مساوی ۰.۳ به اپسيلون مساوی ۰.۰۲ مساوی ۳.۷ درصد است.

%% ============================================================
\section{۳.۷ درصد يعنی چه؟}

يعنی عرض لايه مرزی اثر متوسطی بر لرزش حالت ماندگار برای اين سيستم و اين بهره‌های تنظيم‌شده با \lr{PSO} دارد. بهره سوئيچينگ \lr{K} و اندازه اغتشاش اثر بسيار بزرگ‌تری بر لرزش نسبت به اپسيلون دارند وقتی مسير روی سطح لغزشی يا نزديک آن است.

نقطه نزديک-بهينه تعيين‌شده در مطالعه اپسيلون مساوی ۰.۰۲ است. اين مقدار لرزش را در برابر دقت رديابی قابل قبول متعادل می‌کند. گزارش درباره زبان دقيق است: «نزديک-بهينه» نه «بهينه»، چون جارو فقط سه مقدار اپسيلون را پوشش می‌دهد.

%% ============================================================
\section{چرا نتيجه اصلاح‌شده را باور کنيم؟}

يک مشاور تقريباً حتماً می‌پرسد: چرا ۳.۷ درصد به جای ۶۶.۵ درصد؟

پاسخ دو بخش دارد.

اول، استدلال فنی. روش \lr{RMS} تحت سلطه گذرا شناخته‌شده است که به سوگيری منجر می‌شود. هر معياری که فاز دسترسی را از حالت ماندگار جدا نکند دو پديده متفاوت را مخلوط می‌کند. روش حوزه فرکانسی با طراحی از اين اجتناب می‌کند.

دوم، استدلال فيزيکی. عدد ۶۶.۵ درصد نشان می‌داد که تغيير اپسيلون با ضريب ۱۵ \dash{} از ۰.۳ به ۰.۰۲ \dash{} تقريباً لرزش را حذف می‌کند. اين برای يک کنترل‌گر که در آن بهره سوئيچينگ \lr{K} مساوی ۲.۲۳ است و کران اغتشاش غيرصفر است از نظر فيزيکی نامحتمل است. عدد ۳.۷ درصد از نظر فيزيکی با آنچه مکانيزم لايه مرزی بايد مشارکت کند سازگار است.

اصلاح در طول حسابرسی داخلی گزارش‌های \lr{MT-5} و \lr{MT-6} به دست آمد. توسط شخص ثالث مستقل تأييد نشده. اين يک محدوديت مشروع برای اعتراف صادقانه است.

%% ============================================================

\begin{mdframed}[
  backgroundcolor=audiotan,
  linecolor=audionote!70,
  linewidth=1pt,
  innerleftmargin=12pt,
  innerrightmargin=12pt,
  innertopmargin=8pt,
  innerbottommargin=8pt
]
\textbf{\color{audionote}[يادداشت صوتی:]}
در همه شبيه‌سازی‌های عملياتی و اجراهای معيار \dash{} \lr{MT-5}، مونت‌کارلو، جداول عملکرد \dash{} اپسيلون مساوی ۰.۳ است. مطالعه \lr{MT-6} اپسيلون مساوی ۰.۰۲ را به عنوان نزديک-بهينه برای لرزش يافت. چرا اجراهای معيار از ۰.۳ استفاده می‌کنند؟ دو دليل: مطالعه \lr{MT-6} پس از تکميل معيارها نهايی شد، و تغيير به ۰.۰۲ بدون تنظيم مجدد \lr{PSO} نيمرخ عملکرد را تغيير می‌دهد چون بهره‌ها برای اپسيلون ۰.۳ کاليبره شده‌اند. اين يک ناسازگاری برای پنهان کردن نيست. يک انتخاب طراحی مستند با يک اقدام آينده صريح است.
\end{mdframed}

\vspace{0.6cm}

%% ============================================================
\section{نتيجه‌گيری}

داستان اصلاح \lr{MT-6} در واقع شواهد صداقت پژوهشی است. ادعای اصلی به صورت انتقادی بررسی شد، روش‌شناسی ناقص تشخيص داده شد، و نتيجه با روش بهتر و توضيح شفاف اصلاح شد. عدد اصلاح‌شده ۳.۷ درصد از نظر فيزيکی محتمل و از نظر روش‌شناختی سالم است.

آنچه باقی می‌ماند: يک جارو دقيق‌تر اپسيلون، و معيارهای به‌روزشده با استفاده از اپسيلون ۰.۰۲ با بهره‌های تنظيم مجدد‌شده.

در قسمت بعدی: عملکرد سيستم تحت عدم قطعيت پارامتر و اغتشاشات خارجی.

%% ============================================================
\vspace{1.5cm}
\begin{center}
{\color{seriesblue}\rule{0.45\linewidth}{0.4pt}}\\[8pt]
{\small\color{sectiongray}
منابع گزارش: بخش ۶
}
\end{center}

\end{document}
